\newpage
\section{Introdução a computação quântica.}

Nesta seção, buscamos responder às seguintes perguntas: 
\begin{center}
    \textit{O que são qubits e como os computadores quânticos os criam?}
\end{center}
para responder essa perguntar iremos ...

O poder dos computadores quânticos advém da maneira distinta como eles representam e manipulam informações. Computadores quânticos utilizam bits especiais, chamados de bits quânticos, ou qubits. Assim como os bits clássicos, os qubits também correspondem a algo físico. Nosso foco no momento será nos conceitos essenciais da teoria com perspectiva de desenvolvimento de software e algoritmos.
Nos dedicaremos no capítulo  7 sobre o desenvolvimento físico dos qubits.



\subsection{Bit quântico.}


\subsection*{Representação matemática do estado de um qubit.}

\subsection{Normalização dos estados.}

\subsection{Produto interno e base ortogonal.}

\subsection{Superposição.}

\subsection{Probabilidade de resultado de medição.}

\subsection*{Amostragem nas medidas.}

\subsection{Operações em estados de um qubits.}

\subsection*{Aplicando uma operação quântica.}

\subsection{Simples algoritmo quântico.}
% \begin{itemize}
%     \item Um meio de medir um qubit para determinar em que estado ele se encontra;
%     \item Uma forma de manipular o estado de um qubit para realizar cálculos.
% \end{itemize}


